\documentclass[11pt,letterpaper]{article}
\usepackage[margin=0.86in,headheight=15pt]{geometry}
\usepackage[T1]{fontenc}
\usepackage[utf8]{inputenc}
\usepackage{lmodern,microtype}
\usepackage{amsmath,amssymb,amsthm,mathtools}
\usepackage{booktabs,longtable,array,enumitem}
\usepackage[dvipsnames]{xcolor}
\usepackage{fancyhdr,hyperref,listings}
\definecolor{ink}{HTML}{17344A}
\definecolor{soft}{HTML}{EDF2F5}
\hypersetup{colorlinks=true,linkcolor=ink,citecolor=ink,urlcolor=ink,
 pdftitle={AC-02: Exact certificates and a component obstruction},
 pdfauthor={ChatGPT research report},pdfsubject={Partial results on complex rank of 3 by 3 matrix multiplication}}
\pagestyle{fancy}\fancyhf{}
\lhead{\small\textcolor{ink}{AC-02 \;|\; Exact certificates}}
\rhead{\small Research report}
\cfoot{\small\thepage}
\renewcommand{\headrulewidth}{0.3pt}
\setlength{\parskip}{5pt}\setlength{\parindent}{0pt}
\setlength{\emergencystretch}{3em}
\newtheorem{theorem}{Theorem}[section]
\newtheorem{lemma}[theorem]{Lemma}
\newtheorem{proposition}[theorem]{Proposition}
\newtheorem{corollary}[theorem]{Corollary}
\theoremstyle{definition}\newtheorem{definition}[theorem]{Definition}
\theoremstyle{remark}\newtheorem{remark}[theorem]{Remark}
\newcommand{\CC}{\mathbb C}\newcommand{\ZZ}{\mathbb Z}\newcommand{\QQ}{\mathbb Q}
\newcommand{\Cstar}{\mathbb C^{\!*}}
\DeclareMathOperator{\rank}{rank}\DeclareMathOperator{\tr}{tr}
\DeclareMathOperator{\GL}{GL}\DeclareMathOperator{\diag}{diag}
\newcommand{\Span}{\operatorname{span}}
\newcommand{\file}[1]{\nolinkurl{#1}}
\lstset{basicstyle=\small\ttfamily,breaklines=true,columns=fullflexible,frame=single,
 rulecolor=\color{soft},backgroundcolor=\color{soft},xleftmargin=5pt,xrightmargin=5pt}
\setcounter{tocdepth}{1}
\begin{document}
\begin{titlepage}
\vspace*{0.35in}
{\large\sffamily\color{ink} OPEN PROBLEMS IN NUMERICAL LINEAR ALGEBRA}\par
\vspace{0.55in}
{\Huge\bfseries\color{ink} AC-02}\par\vspace{0.2in}
{\LARGE\bfseries Exact certificates and a\par component obstruction}\par
\vspace{0.25in}
{\large Complex bilinear rank of the $3\times3$ matrix product}\par
\vspace{0.35in}
{\large Research report and reproducible verification package}\par
14 September 2026\par
\vspace{0.4in}
\begingroup\setlength{\fboxsep}{13pt}
\colorbox{soft}{\begin{minipage}{0.91\textwidth}
\textbf{Status of the original problem: not resolved by this report.}\par
The exact minimum is not established. The verified general interval remains
$19\le R_{\CC}(M_3)\le23$. No rank-22 identity and no global lower bound of 23
are supplied. The results below concern precisely specified families and
necessary conditions for arbitrary decompositions.
\end{minipage}}\endgroup\par
\vspace{0.3in}
\textbf{Main computer-assisted result.}
The complete coefficient variety with the displayed Laderman zero pattern
is an algebraic torus of dimension 58. Its closure under the standard three
matrix changes of basis is a 76-dimensional irreducible component of the
23-slot Brent variety. Every slot has a fixed nonzero polynomial signature
on this component, so no zero-slot point lies on it.

\textbf{Additional results.}
Exact support obstructions are verified for Laderman's and Sun's
23-product schemes. A reshuffling argument gives general matrix-rank-sum
constraints and a sharper classification of the output-rank profile in a
tight hypothetical 22-product decomposition.

\vfill
\small
The archive contains the proofs, exact integer and Laurent certificates,
a replay checker, attributed algorithm data, rational family instances,
and numerical exploration logs. The checking is not a Lean formalization
and has not undergone an independent external mathematical audit.
No claim of priority for the restricted results is made.
\end{titlepage}
\tableofcontents
\newpage

\section{The target and the result actually obtained}
\label{sec:scope}
AC-02 asks for the least number of separated bilinear products needed to
multiply arbitrary complex $3\times3$ matrices. Each product must multiply
one linear form in the first input by one linear form in the second input;
scalar linear combinations are free. This is the exact tensor-rank model,
not an addition-count, approximate, border-rank, or mixed-input model
\cite{ac02}.

The established bounds used here are
\begin{equation}\label{eq:interval}
19\le R_{\CC}(M_3)\le23.
\end{equation}
Bl\"aser's survey records these bounds in its introduction and Example 5.10
\cite{blaser}. The lower bound is imported from the literature, not reproved
by this report. The upper bound is verified here from an explicit Laderman
identity, and separately from Sun's displayed identity \cite{laderman,sun}.
The source checks did not locate a qualifying improvement for the exact
complex model; the repository's current statement and Alman--Li also
identify the exact-rank question as unresolved \cite{ac02,almanli}.

Wang's finite-field lower-bound work has a different field of coefficients
\cite{wang}. In particular, a nonexistence result for coefficients in
$\mathbb F_2$ does not establish nonexistence over $\CC$. No finite-field
search result is used here to infer a complex tensor-rank lower bound.

\subsection{An exact separation of conclusions}
The principal conclusions proved with the accompanying finite certificates
are summarized below. A ``slot'' means one ordered rank-one tensor summand,
including its three coefficient vectors.

\begin{center}
\begin{tabular}{p{0.57\textwidth}p{0.33\textwidth}}
\toprule
Conclusion & Scope \\
\midrule
All 729 identities hold for two 23-slot decompositions & Existing upper bounds \\
No factor vector can vanish with the prescribed supports & Laderman and Sun supports \\
All three pairwise Khatri--Rao matrices have column rank 23 & The same support varieties \\
The complete support variety is $(\Cstar)^{58}$ & Laderman support \\
A 76-dimensional component contains no zero-slot point & The component through Laderman \\
$\sum_t\rank W_t\ge27$, and analogues for $U,V$ & Every exact decomposition \\
A tight 22-slot output profile is $2^5 1^{17}$ & Conditional on $\sum_t\rank W_t=27$ \\
\bottomrule
\end{tabular}
\end{center}

These conclusions neither exhibit a 22-product algorithm nor exclude all
such algorithms. In particular, the component theorem does not assert that
every decomposition is equivalent to Laderman's. The failure of that global
coverage statement is the central gap, not a missing computational check
within the displayed component.

\section{Coordinates, identities, and the coefficient variety}
\label{sec:coordinates}
Throughout, matrix indices $i,j,k$ range over $\{0,1,2\}$. A matrix entry
$(i,j)$ has vector coordinate $3i+j$. Product slots are indexed
$t=1,\ldots,r$ in the text and $t=0,\ldots,r-1$ in JSON and Python unless
explicitly labeled otherwise.

Write $U_t,V_t,W_t$ for $3\times3$ coefficient matrices and define
\begin{equation}\label{eq:algorithm}
AB=\sum_{t=1}^r
 \left(\sum_{i,j}U_{t,ij}A_{ij}\right)
 \left(\sum_{j,k}V_{t,jk}B_{jk}\right)W_t.
\end{equation}
There is no complex conjugation in this formula or in any algebraic
signature below. Equality for every input pair is equivalent to the 729
Brent equations
\begin{equation}\label{eq:brent}
\sum_{t=1}^r U_{t,ij}V_{t,j'k}W_{t,i'k'}
  =\delta_{ii'}\delta_{jj'}\delta_{kk'}.
\end{equation}
Let $T\in\ZZ^{9\times9\times9}$ have entry 1 at
$(3i+j,3j+k,3i+k)$ and 0 elsewhere. It has 27 entries equal to 1 and 702
equal to 0. Equation \eqref{eq:brent} is
$T=\sum_t u_t\otimes v_t\otimes w_t$, with row-major vectorized factors.

\begin{definition}
$X_{23}$ is the affine variety in $\CC^{621}$ cut out by
\eqref{eq:brent} with $r=23$. All slots and all factors are ordered.
Zero factor vectors are allowed. For a fixed baseline $F^{(0)}=(U,V,W)$,
its support variety consists of the points in $X_{23}$ at which every
coordinate zero in $F^{(0)}$ remains zero. Permitted coordinates are not
assumed nonzero in this definition.
\end{definition}

A point has a zero slot exactly when at least one of its three vectors in
that slot is zero: a tensor product of three nonzero vectors over a field
is nonzero. Hence
\begin{equation}\label{eq:zeroequiv}
R_{\CC}(M_3)\le22
\quad\Longleftrightarrow\quad
X_{23}\text{ contains a zero-slot point}.
\end{equation}
The forward implication pads an algorithm with zeros; the reverse
implication deletes the zero slot. An algorithm of fewer than 22 products
can first be padded to 22 products as well.

\subsection{Data and baseline verification}
The files \file{data/laderman23.json} and \file{data/sun23.json} contain three
$23\times9$ signed-integer arrays. They are reconstructed from explicit
scalar formulas by \file{code/build_data.py}. Laderman's formulas are also
printed in Appendix \ref{app:laderman}. Sun's data are transcribed from
Section 5 of his paper \cite{sun}; the data are not an algorithm discovered
in this work.

The checker evaluates every left side of \eqref{eq:brent} as an integer and
compares it to the corresponding Kronecker delta. All 729 checks pass for
each baseline. This proves exact identities, not merely agreement on
random numerical inputs. Because the identities have integer coefficients,
the displayed algorithms also remain identities under extension to $\CC$.
Their optimality does not follow from verification.

\section{Nonvanishing on two prescribed support varieties}
\label{sec:forcing}
For each permitted coordinate, introduce a normalized variable
$x_a=F_a/F_a^{(0)}$. Baseline coefficients are $\pm1$, so the normalization
is harmless even before nonvanishing has been established.
Each restricted Brent equation is a sum of signed cubic monomials, possibly
including the constant $-1$.

\begin{lemma}[Two-monomial propagation]\label{lem:forcing}
Suppose a polynomial equation over a field is
$c_1m_1+c_2m_2=0$, where $c_1,c_2$ are nonzero scalars and $m_1,m_2$ are
monomials. If every variable in $m_1$ is nonzero at every solution of the
system, then every variable in $m_2$ is also nonzero at every solution.
The empty monomial 1 is nonzero without any assumptions.
\end{lemma}
\begin{proof}
The first monomial is nonzero. The equation gives
$m_2=-c_1m_1/c_2\ne0$. A product of field elements is nonzero only if each
factor is nonzero. Applying the argument successively preserves its
universal quantifier over all solutions.
\end{proof}

For example, the coefficient of $A_{00}B_{00}$ in output $(0,0)$, on
Laderman's support, gives
\[
 x_{U_{6,00}}x_{V_{6,00}}x_{W_{6,00}}=1.
\]
These three coordinates are therefore forced nonzero. Other equations
propagate nonvanishing from already established coordinates.

\begin{proposition}[Exact support obstruction]\label{prop:support}
On the Laderman support variety, every one of the 153 permitted
coordinates is nonzero. On the Sun support variety, at least 162 of the
175 permitted coordinates are nonzero. In both varieties all 69 factor
vectors are nonzero, and consequently every one of the 23 slots is nonzero.
\end{proposition}
\begin{proof}
The finite traces in \file{proofs/laderman23_forcing.json} and
\file{proofs/sun23_forcing.json} contain respectively 87 and 99 applications
of Lemma \ref{lem:forcing}. Each step specifies the Brent polynomial, the
monomial already known to be nonzero, and the newly forced variables.
The checker reconstructs the equations from the baseline data; it does
not accept a polynomial supplied by the certificate. It verifies the
hypotheses and conclusion of every step and then checks that each of the
69 factor vectors contains a forced coordinate. The resulting sets have
sizes 153 and 162, respectively.
\end{proof}

This does not assume that the original $\pm1$ entries stay $\pm1$. The
allowed entries may vary arbitrarily in $\CC$, including attempted zeros.
It is precisely the equations that prevent the stated coordinates from
vanishing. For Sun, no claim is made that the remaining 13 permitted
coordinates must be nonzero.

\subsection{A stronger pairwise independence certificate}
For two modes, define the $81\times23$ matrix
\[
 K_{UV}(:,t)=u_t\otimes v_t,
\]
and define $K_{VW}$ and $K_{WU}$ similarly. The ordering of their 81 rows is
lexicographic in the two vector coordinates.

\begin{lemma}[Triangular peeling]\label{lem:peel}
Suppose the columns of a matrix have a fixed permitted support. Assume
there is an ordering of the columns and a row for each column such that,
among the columns not previously removed, the selected row has only that
column as a permitted nonzero, and its entry is known to be nonzero.
Then the columns are linearly independent for every permitted assignment.
\end{lemma}
\begin{proof}
In a linear dependence, the first selected row forces the coefficient of
its unique surviving column to be zero. Remove that column from the
relation and repeat. Every coefficient is zero.
\end{proof}

\begin{proposition}\label{prop:pairwise}
For every solution in either support variety of Proposition
\ref{prop:support}, each of $K_{UV},K_{VW},K_{WU}$ has column rank 23.
\end{proposition}
\begin{proof}
The two \file{*_pairwise.json} certificates each give three lists of
23 row--column choices. The checker verifies the support uniqueness at
each step and verifies that the two coordinates whose product is the
selected entry belong to the forced-nonzero set. Lemma \ref{lem:peel}
applies to each list.
\end{proof}

As one consequence, with the two input-factor arrays fixed at any such
solution, the output-factor array is uniquely determined by the tensor
identity: the matrix equation is $K_{UV}W=T_{UV|W}$ and $K_{UV}$ is
injective. Reoptimizing only the output factors therefore cannot remove a
slot. This is an obstruction for the displayed pairs, not a lower bound
against all alternative pairs of factors.

\section{The complete Laderman support variety is a torus}
\label{sec:torus}
Let $S\subseteq X_{23}$ denote the Laderman support variety. Proposition
\ref{prop:support} places all its 153 allowed coordinates in $\Cstar$.
Thus its equations may be treated as Laurent polynomial equations in the
normalized coordinates $x_1,\ldots,x_{153}$.

\subsection{Reduction to unit-monomial equations}
Of the 729 restricted equations, 610 vanish identically. The other 119
have the following forms: 23 equations have one cubic and a nonzero
constant; 92 have two cubics and no constant; and four have three cubics
and a constant. The first 115 therefore become equations $x^a=1$, with
integer exponent vectors $a$.

Let $A_0\in\ZZ^{115\times153}$ consist of these exponent rows. The
certificate establishes $\rank A_0=91$ and gives a $91\times91$ minor
whose inverse is an integer matrix. It also gives
$E_0\in\ZZ^{153\times62}$ with an identity block on 62 free coordinates
and $A_0E_0=0$.

In each of the four remaining equations, two oppositely signed cubics
have exponent difference in the integer row lattice generated by the
selected rows of $A_0$. Explicit integer combinations in the certificate
prove these four lattice memberships. Therefore those two cubics agree
on every solution of $x^{A_0}=1$ and cancel. The survivor gives one further
equation $x^a=1$.

For clarity, the four affected coefficient coordinates are listed below.
The entries in the first column are 0-based vector coordinates $(a,b,c)$;
slot numbers in the last column are 1-based. All factor entries in that
column are normalized by their baseline signs.
\begin{center}
\begin{tabular}{ccp{0.47\textwidth}}
\toprule
$(a,b,c)$ & Brent index $81a+9b+c$ & Surviving monomial \\
\midrule
$(3,1,4)$ & 256 & $x_{U_{5,10}}x_{V_{5,01}}x_{W_{5,11}}$ \\
$(5,8,5)$ & 482 & $x_{U_{18,12}}x_{V_{18,22}}x_{W_{18,12}}$ \\
$(6,2,8)$ & 512 & $x_{U_{9,20}}x_{V_{9,02}}x_{W_{9,22}}$ \\
$(8,7,7)$ & 718 & $x_{U_{15,22}}x_{V_{15,21}}x_{W_{15,21}}$ \\
\bottomrule
\end{tabular}
\end{center}
Let $A\in\ZZ^{119\times153}$ be $A_0$ with these four exponent rows
appended. Its certified rank is 95, and a selected $95\times95$ minor has
an integer inverse.

\begin{lemma}[Unimodular monomial elimination]\label{lem:unimod}
Let $A\in\ZZ^{m\times n}$. Select $r$ rows and $r$ pivot columns so that
the resulting square matrix $B$ has an integer inverse. Let $C$ be the
selected rows restricted to the other $n-r$ columns. Define
$E\in\ZZ^{n\times(n-r)}$ by
\[
 E_{\mathrm{free},:}=I_{n-r},\qquad
 E_{\mathrm{pivot},:}=-B^{-1}C.
\]
If $AE=0$, then the solution set of $x^A=1$ in $(\Cstar)^n$ is
isomorphic to $(\Cstar)^{n-r}$ by $x_a=\prod_j z_j^{E_{aj}}$.
\end{lemma}
\begin{proof}
For every $z$, the equations hold because their exponent vectors become
zero after multiplication by $E$. Conversely, let $z$ be the free
coordinates of a solution. The selected equations are
$x_{\mathrm{pivot}}^Bz^C=1$. Taking the integer monomial combinations
specified by $B^{-1}$ expresses each pivot variable uniquely as the
corresponding row of $-B^{-1}C$. No extraction of roots is involved.
All other equations hold for the parametrization by $AE=0$.
Both the map and its inverse, coordinate projection to $z$, are regular
maps of tori.
\end{proof}

\begin{theorem}[Complete fixed-support parametrization]\label{thm:torus}
There is an explicit isomorphism
\[
 S\simeq (\Cstar)^{58}.
\]
More precisely, let $E$ be the $153\times58$ integer matrix in
\file{proofs/laderman_torus.json}, under \file{final.E}. Then every point
of $S$, and only such a point, is
\begin{equation}\label{eq:param}
 F_a(z)=F_a^{(0)}\prod_{j=1}^{58}z_j^{E_{aj}},\qquad
 z\in(\Cstar)^{58},
\end{equation}
with all other coordinates zero. Every entry of $E$ has absolute value
at most 2.
\end{theorem}
\begin{proof}
Nonvanishing follows from Proposition \ref{prop:support}. The preceding
integer row-lattice certificates reduce the restricted equations to
$x^A=1$. The certificate contains the selected rows and pivot columns,
the integer inverse, the 58 free indices, and $E$. The replay checker
verifies both inverse identities, the identity block on the free
coordinates, the pivot formula in Lemma \ref{lem:unimod}, and $AE=0$.
That lemma proves completeness and uniqueness.

As a redundant check of sufficiency, the checker substitutes
\eqref{eq:param} into each of the 729 original equations, groups terms
by their exact 58-dimensional exponent vectors, and verifies that every
resulting integer coefficient is zero. Thus no ignored residual equation
or numerical approximation enters the parametrization.
\end{proof}

The integer inverse is important. A rational rank computation alone would
not exclude additional finite components, such as those arising from
$x^2=1$. Here the selected minors are unimodular, so those extra root
choices do not occur.

\subsection{A six-parameter slice after elementary rescalings}
There are 46 independent term-rescaling parameters:
\[
 U_t\mapsto\lambda_tU_t,\quad V_t\mapsto\mu_tV_t,\quad
 W_t\mapsto(\lambda_t\mu_t)^{-1}W_t.
\]
Also allow diagonal basis changes
$P=\diag(p_1,p_2,1)$, $Q=\diag(q_1,q_2,1)$,
$R=\diag(r_1,r_2,1)$ in the action of Section \ref{sec:component}.
These give a 52-parameter torus of support-preserving transformations.

\begin{proposition}[Explicit elementary-gauge slice]\label{prop:slice}
Every point of $S$ has a unique representative, under this 52-parameter
rescaling group, for which 52 designated free coordinates in
\eqref{eq:param} equal 1. The other six free coordinates have 0-based
indices
\[
21,\ 37,\ 45,\ 47,\ 50,\ 52.
\]
The representative is obtained by retaining precisely these six columns
of $E$ and setting all other parameters to 1.
\end{proposition}
\begin{proof}
Let $B$ be the $153\times55$ exponent matrix for all 46 term rescalings
and all nine diagonal entries of $P,Q,R$. Let $D$ be its restriction to
the 58 free coordinate rows. Then $ED=B$; this identity is checked over
$\ZZ$.

Select the 52 columns corresponding to all term rescalings and the first
two diagonal entries of each matrix. The certificate gives 52 rows of
$D$ whose square submatrix has an integer inverse. These rows are exactly
the complement of the six displayed indices. For any free-coordinate
vector $z$, the monomial equations making these 52 coordinates equal to
1 have a unique solution for the 52 rescaling parameters, by the same
integer-inverse argument as Lemma \ref{lem:unimod}.
The remaining three diagonal generators are integer monomial
combinations of the selected generators on $S$, as is also verified in
the certificate.
\end{proof}

This is a quotient by the stated elementary rescalings, not a claim that
the six coordinates classify algorithms under every possible equivalence.
The file \file{data/laderman_six_parameter_exponents.csv} prints all 153
coefficient formulas as baseline signs and six exponent columns. The
script \file{code/instantiate_family.py} creates exact rational instances.
The sample with parameters $(2,3,5,7,11,13)$ passes all 729 identities.

Parameterized matrix multiplication algorithms and heuristic ways of
introducing parameters have substantial prior literature. In particular,
Heule--Kauers--Seidl report that their particular iterative linear
parameter-introduction procedure found no parameters for Laderman's
scheme \cite[Section 7]{hks}. That is a statement about the output of a
specified heuristic, not a theorem that every coupled nonlinear
coefficient deformation is trivial. No priority claim is made here for
the present torus calculation.

\section{A 76-dimensional component without zero slots}
\label{sec:component}
The support theorem is not confined to algorithms that keep their zero
patterns after a basis change. We now take an orbit closure inside the
full coefficient variety $X_{23}$.

\subsection{The action and a polynomial invariant}
Let $G=\GL_3(\CC)^3$. Its action on coefficient matrices is
\begin{equation}\label{eq:action}
\begin{split}
 U'_t&=P^\top U_tQ^{-\top},\\
 V'_t&=Q^\top V_tR^{-\top},\\
 W'_t&=P^{-1}W_tR.
\end{split}
\end{equation}
It preserves \eqref{eq:algorithm}: apply the original algorithm to
$A'=PAQ^{-1}$ and $B'=QBR^{-1}$, and transform its output by
$P^{-1}(\,\cdot\,)R$. In particular, it preserves $X_{23}$.

For a slot define the polynomial
\begin{equation}\label{eq:sig}
 \sigma_t(F)=\tr(U_tV_tW_t^\top)
 =\sum_{i,j,k}U_{t,ij}V_{t,jk}W_{t,ik}.
\end{equation}
This is the familiar contraction signature in the theory of matrix
multiplication decompositions \cite{tichavsky}. Its invariance can be
checked directly here:
\[
 U'_tV'_t(W'_t)^\top
   =P^\top U_tV_tW_t^\top P^{-\top}.
\]
Taking trace proves $\sigma_t(F')=\sigma_t(F)$.

\begin{proposition}[Constant signatures on the torus]\label{prop:sig}
On $S$ the signatures are
\begin{equation}\label{eq:sigvalues}
\sigma_t=\begin{cases}
2,&t\in\{4,7,12,16\},\\
1,&\text{all other slots}.
\end{cases}
\end{equation}
The slot numbering is that of the included Laderman data.
\end{proposition}
\begin{proof}
Substitute \eqref{eq:param} in \eqref{eq:sig}. The checker groups the
resulting Laurent monomials by exponent vector. Every nonconstant
coefficient is zero, and the constant coefficients are exactly
\eqref{eq:sigvalues}. This holds symbolically for all 58 nonzero complex
parameters, not just at the baseline point.
\end{proof}

Define
\begin{equation}\label{eq:Z}
 Z=\overline{G\cdot S}^{\,\mathrm{Zar}}\ \subseteq X_{23},
\end{equation}
where the closure is in affine coefficient space.

\begin{theorem}[Orbit-closure obstruction]\label{thm:closure}
$Z$ is irreducible, and no point of $Z$ has a zero slot.
\end{theorem}
\begin{proof}
Both $G$ and $S$ are irreducible, so the closure of the image of the
regular action map $G\times S\to X_{23}$ is irreducible. Every polynomial
$\sigma_t-c_t$, with $c_t$ from \eqref{eq:sigvalues}, vanishes on $G\cdot S$
and hence on its Zariski closure. At a zero slot, its cubic signature
would be zero. This contradicts $c_t\in\{1,2\}$.
\end{proof}

This theorem also covers all finite coordinatewise limits of transformed
family members. It does not cover limits in which the coefficient arrays
diverge while only their summed tensor converges.

\subsection{Exact local dimension and identification of the component}
Let $L=F^{(0)}$ be the baseline Laderman point. Let $J$ be the Jacobian of
the 729 Brent polynomials at $L$, with columns ordered by mode, slot,
then row-major entry. It is an integer $729\times621$ matrix.
For example, the column for the variable $U_{t,a}$ is
$e_a\otimes v_t\otimes w_t$, vectorized in lexicographic order.

Differentiate \eqref{eq:param} at $z=(1,\ldots,1)$. Its 58 derivative
columns have entry $F_a^{(0)}E_{aj}$ at each allowed coordinate and zero
at the forbidden coordinates. Add 18 derivatives of \eqref{eq:action},
one for each off-diagonal elementary matrix $D$ in each group factor:
\begin{equation}\label{eq:directions}
\begin{array}{c|ccc}
 &\delta U_t&\delta V_t&\delta W_t\\ \hline
P&D^\top U_t&0&-DW_t\\
Q&-U_tD^\top&D^\top V_t&0\\
R&0&-V_tD^\top&W_tD
\end{array}
\end{equation}
All derivatives are taken at the identity. Denote the resulting
$621\times76$ matrix by $H$.

\begin{lemma}[Exact tangent certificates]\label{lem:tangent}
Over $\CC$, $\rank H=76$ and $\rank J=545$. Also $JH=0$ over $\ZZ$.
\end{lemma}
\begin{proof}
The checker reconstructs both matrices from the baseline and the
monomial exponent matrix. It verifies $JH=0$ with unbounded integer
arithmetic. The file \file{proofs/laderman_component.json} lists a
$76\times76$ minor of $H$ and a $545\times545$ minor of $J$ that are
nonsingular modulo the prime $p=1000003$. Primality is checked by trial
division, and nonsingularity is checked by exact modular elimination.

An integer determinant nonzero modulo $p$ is a nonzero integer, hence is
nonzero over $\CC$. Thus $\rank H\ge76$ and $\rank J\ge545$. The size of
$H$ and the relation $JH=0$ imply
$\rank J\le621-76=545$. All claimed ranks follow.
\end{proof}

\begin{theorem}[The Laderman component]\label{thm:component}
$Z$ is the unique irreducible component of $X_{23}$ through $L$.
It has dimension 76, and $L$ is a nonsingular point of $X_{23}$.
In particular, this entire irreducible component has no zero-slot point.
\end{theorem}
\begin{proof}
The columns of $H$ are derivatives of the action map $G\times S\to X_{23}$
at $(I,L)$. Lemma \ref{lem:tangent} gives differential rank 76, so its
image closure $Z$ has dimension at least 76. For instance, a nonzero
76-dimensional derivative minor yields local analytic image dimension
at least 76 over $\CC$.

Every tangent vector to $X_{23}$ at $L$ lies in $\ker J$. Therefore
\[
\dim_L X_{23}\le\dim T_LX_{23}\le621-545=76.
\]
Since $Z$ is irreducible and contains $L$, its dimension gives the reverse
lower bound $\dim_LX_{23}\ge\dim Z\ge76$. Consequently all these dimensions
are 76. Equality between local dimension and tangent dimension makes $L$
nonsingular, so exactly one irreducible component of $X_{23}$ passes
through $L$. That component contains $Z$, and the equal dimensions force
it to be $Z$ itself. Theorem \ref{thm:closure} gives the last assertion.
\end{proof}

Jacobian methods and orbit-dimension bounds for Brent varieties appear in
Li--Bao--Zhang \cite{lbz}. The argument here does not assume that the ideal
generated by the Brent polynomials is radical: even if it were not,
$\ker J$ would still be an upper bound for the tangent space of its zero
set. The 76-dimensional family provides the independent lower bound
needed to close the dimension calculation.

\subsection{Precisely what the component theorem excludes}
It excludes reaching a zero slot by any sequence that stays inside $Z$,
including arbitrary changes of basis and finite limiting operations.

It does not show that $X_{23}=Z$. It does not exclude another irreducible
component containing a padded rank-22 algorithm. Nor does it preclude a
continuous path that leaves $Z$ through an intersection with another component
elsewhere. Finally, all 23 slots being nonzero is not itself proof that
the tensor cannot be expressed by a different collection of 22 nonzero
slots. These distinctions prevent the local result from being mistaken
for $R_{\CC}(M_3)=23$.

\section{General reshuffling constraints and the tight case}
\label{sec:reshuffle}
The next arguments apply to arbitrary exact decompositions, without
prescribing support, signs, reality, or a component. Here $n$ may be
arbitrary, and all retained slots are assumed nonzero.

\begin{theorem}[Three matrix-rank-sum inequalities]\label{thm:ranksum}
For an exact decomposition of $n\times n$ matrix multiplication,
\begin{equation}\label{eq:ranksum}
\sum_t\rank U_t\ge n^3,\qquad
\sum_t\rank V_t\ge n^3,\qquad
\sum_t\rank W_t\ge n^3.
\end{equation}
These are ordinary ranks of the coefficient matrices, not tensor ranks.
\end{theorem}
\begin{proof}
Reshape \eqref{eq:brent} into a matrix whose row is $(i,j,k')$ and column
is $(i',j',k)$. The target becomes $I_{n^3}$. A summand becomes
\[
 K_t[(i,j,k'),(i',j',k)]=U_{t,ij}V_{t,j'k}W_{t,i'k'}.
\]
Factor $W_t=X_tY_t^\top$ with $q_t=\rank W_t$ independent columns in each
of $X_t,Y_t$. Then $K_t=L_tR_t$, where
\[
 L_t[(i,j,k'),a]=U_{t,ij}Y_t[k',a],\qquad
 R_t[a,(i',j',k)]=X_t[i',a]V_{t,j'k}.
\]
Thus $\rank K_t\le q_t$ (indeed equality holds for nonzero $U_t,V_t$).
Subadditivity of matrix rank applied to $I_{n^3}=\sum_tK_t$ gives the
$W$ inequality. Cyclically permuting the three factors in the trace
form of the multiplication tensor gives the other two inequalities;
transposes do not change ordinary matrix ranks.
\end{proof}

For a 22-slot decomposition with $n=3$, each inequality says
$\sum_t(\rank W_t-1)\ge5$, and similarly for $U$ and $V$.
This alone is much weaker than a lower bound of 23 products.

\begin{definition}
A decomposition is \emph{output-tight} when
$\sum_t q_t=n^3$, where $q_t=\rank W_t$.
\end{definition}

\begin{theorem}[Dual block identities in the tight case]\label{thm:tight}
For an output-tight decomposition, choose any full-column-rank
factorizations $W_t=X_tY_t^\top$. Then for all slots $s,t$,
\begin{equation}\label{eq:blocks}
X_s^\top U_tV_sY_t=
\begin{cases}
I_{q_t},&s=t,\\
0_{q_s\times q_t},&s\ne t.
\end{cases}
\end{equation}
In particular,
\begin{equation}\label{eq:tightself}
q_t\le\min\{\rank U_t,\rank V_t\},\qquad
\sigma_t=q_t.
\end{equation}
\end{theorem}
\begin{proof}
Concatenate the matrices in the proof of Theorem \ref{thm:ranksum}:
$L=[L_1\ \cdots\ L_r]$ and
$R=[R_1^\top\ \cdots\ R_r^\top]^\top$.
The exact identity gives $LR=I_{n^3}$.
Tightness makes $L,R$ square, hence $RL=I_{n^3}$ as well.
The $(s,t)$ block of $RL$ is
\[
\sum_{i,j,k}X_s[i,b]V_s[j,k]U_t[i,j]Y_t[k,a]
   =\bigl(X_s^\top U_tV_sY_t\bigr)_{ba},
\]
proving \eqref{eq:blocks}. The diagonal blocks imply the rank inequalities
in \eqref{eq:tightself}. Their traces give
$q_t=\tr(X_t^\top U_tV_tY_t)
=\tr(U_tV_tW_t^\top)=\sigma_t$.
\end{proof}

\begin{theorem}[An invertible output forces 25 slots]\label{thm:fullrank}
For $n=3$, an output-tight decomposition with a rank-three output matrix
has exactly 25 nonzero slots. Thus an output-tight decomposition with at
most 24 nonzero slots has no rank-three output matrix.
\end{theorem}
\begin{proof}
Let $q_0=3$. Then $X_0,Y_0$ are invertible, and the self-block identity
$X_0^\top U_0V_0Y_0=I_3$ makes $U_0,V_0$ invertible.
For any other slot $t$, the off-diagonal block with $s=0$ yields
\[
U_tV_0Y_t=0.
\]
The matrix $V_0Y_t$ has column rank $q_t$, so
$\rank U_t\le3-q_t$.
Using the block with second slot index 0 similarly gives
\[
X_t^\top U_0V_t=0,\qquad\rank V_t\le3-q_t.
\]
By \eqref{eq:tightself}, $q_t\le\rank U_t,\rank V_t$.
Thus $2q_t\le3$, and every other $q_t=1$.
Output tightness is now $27=3+(r-1)$, hence $r=25$.
\end{proof}

\begin{corollary}[The output-tight 22-slot profile]\label{cor:22}
An output-tight 22-slot decomposition of $M_3$ must have exactly five
rank-two outputs and seventeen rank-one outputs. It has no rank-three
output. Its contraction signatures have the corresponding multiset
$\{2,2,2,2,2,1,\ldots,1\}$.
\end{corollary}
\begin{proof}
By Theorem \ref{thm:fullrank}, the output ranks lie in $\{1,2\}$. If $a$
of them equal 2, output tightness gives $27=2a+(22-a)$, so $a=5$.
The signature claim follows from Theorem \ref{thm:tight}.
\end{proof}

No argument here establishes that an arbitrary rank-22 decomposition must
be output-tight in any orientation. Therefore Corollary \ref{cor:22}
does not exclude every possible output profile.

\subsection{Application to the Sun baseline}
For the supplied Sun data, direct exact matrix-rank computations give
\[
 \sum_t\rank U_t=32,\qquad
 \sum_t\rank V_t=30,\qquad
 \sum_t\rank W_t=27.
\]
The four rank-two outputs are slots $5,8,10,23$; all other outputs have
rank one. The certificate \file{proofs/sun_tight.json} gives explicit
integer factorizations $W_t=X_tY_t^\top$. The checker verifies their
ranks, forms the $27\times27$ matrices $L,R$, checks both $LR=I$ and
$RL=I$, and independently verifies all 529 block identities.

Consequently, one cannot delete any Sun slot while requiring the ordinary
matrix rank of each remaining output to stay at or below its current
value. Such a deletion would lower the output-rank sum below 27, in
contradiction to Theorem \ref{thm:ranksum}, regardless of how the two input
factor arrays are changed. A successful reduction would have to escape
that output-rank restriction.

\section{Verification, reproducibility, and numerical experiments}
\label{sec:verification}
\subsection{Trust boundaries of the certificate checker}
The exact checker does not import the certificate generator and does not
use a symbolic algebra system or floating-point tolerances. It does share
indexing and matrix-building helpers with the generator. Consequently it
is a replay checker, not an independent authorship or formal-verification
claim. A second, stand-alone checker evaluates arbitrary rational-complex
identities directly from six matrix indices and uses only the Python
standard library.

The exact operations used are integer arithmetic, integer matrix products,
Laurent exponent addition, rational elimination for small coefficient
matrices, and finite-field elimination for selected integer minors. Large
integer products use unbounded Python integers. Modular operations use
64-bit integer arrays with prime $1000003$; every multiplication before
reduction is below $10^{12}$, within the integer range.

The modular-minor argument is different from transferring a tensor-rank
lower bound across fields. A specific nonzero integer determinant remains
nonzero over $\CC$. A statement that a polynomial system has no solutions
with coordinates in $\mathbb F_2$ says nothing comparable about its
complex solutions.

\subsection{Commands}
From the root of the extracted archive, with the recorded dependencies
installed:
\begin{lstlisting}
python code/verify.py
python code/self_tests.py
python code/verify_candidate.py data/laderman23.json
python code/verify_candidate.py data/sun23.json
\end{lstlisting}
Regeneration, which additionally uses SymPy, is:
\begin{lstlisting}
python code/build_data.py
python code/generate_certificates.py
python code/verify.py
\end{lstlisting}
A nontrivial exact rational family member can be created by:
\begin{lstlisting}
python code/instantiate_family.py --parameters 2 3 5 7 11 13 \
  --output results/example.json
python code/verify_candidate.py results/example.json
\end{lstlisting}
The candidate checker supports $\QQ(i)$ coefficients represented by exact
rational strings or separate real and imaginary rational parts. It rejects
floating-point coefficients. It does not implement arbitrary algebraic
number fields, so its accepted coefficient format must not be mistaken
for a restriction in the mathematical statement of AC-02.

The seven regression tests include correct real and complex-gauged
identities, a nontrivial rational family member, and deliberately
corrupted identities, forcing traces, and Laurent certificates. All seven
pass. A corruption being rejected is a software regression test, not a
replacement for the proof of the verification method.

\subsection{Unrestricted numerical exploration}
The numerical search frees every coefficient in 22 slots, including
entries that were zero in the starting schemes. It is not a support-fixed
search. The recorded run uses seed 20260914 and 54 starting points:
46 obtained by deleting each slot of each baseline, plus four random real
and four random complex starts. Each receives 500 alternating
least-squares sweeps. The four best starts and two complex starts receive
up to 30 damped Gauss--Newton steps.

The best recorded point, from deleting Sun slot 21, has
\[
 \left\|\sum_{t=1}^{22}u_t\otimes v_t\otimes w_t-T\right\|_F
 =0.16307567682559754,
\]
with relative residual $0.031383928637846385$. This is not an exact
identity. The file \file{results/numerical_rank22_search.json} contains
all settings and iteration records; the best arrays are retained in
\file{results/numerical_best_rank22.npz}.

\textbf{No mathematical lower bound follows from these failures.}
Nor would a very small residual by itself establish an exact rank-22
identity. Coefficients escaping to infinity can represent an approximation
phenomenon rather than a finite exact decomposition. The numerical run is
reported solely as reproducible exploration and is not used in any proof.

\section{The remaining global step}
\label{sec:gap}
Let
\[
 \mathcal N=\{F\in X_{23}:\text{at least one slot of }F\text{ is zero}\}.
\]
By \eqref{eq:zeroequiv}, a proof that $\mathcal N$ is empty would, together
with the verified upper bound, establish $R_{\CC}(M_3)=23$.
Theorem \ref{thm:component} proves only $Z\cap\mathcal N=\varnothing$.
No certificate in this archive establishes emptiness on every other
component of $X_{23}$.

There is an equivalent unrestricted polynomial target. Let
$I_{22}\subset\QQ[U,V,W]$ be the ideal of the 729 Brent polynomials in
594 variables for 22 slots. A global nonexistence proof can be represented
by an exact polynomial identity
\begin{equation}\label{eq:nullstellensatz}
 1=\sum_{a,b,c=0}^{8}h_{abc}(U,V,W)
 \left(\sum_{t=1}^{22}U_{t,a}V_{t,b}W_{t,c}-T_{abc}\right).
\end{equation}
Such an identity would be a directly checkable certificate of
$R_{\CC}(M_3)>22$. Conversely, if the complex zero set is empty, the
Nullstellensatz provides an identity over $\CC$; since the generators
are rational and only finitely many monomial coefficients occur, linear
algebra over $\QQ$ yields rational $h_{abc}$ as well. No identity of the
form \eqref{eq:nullstellensatz} is supplied here.

An exact 22-slot construction is the alternative global direction. It
would improve the upper bound but would still require a matching lower
bound to determine the exact minimum. Restricting such a construction to
small integer coefficients, a prescribed support, a symmetric ansatz, or
a known component would need an additional completeness theorem before
a negative search could resolve AC-02.

\subsection{Consequences for further work}
The component theorem identifies a concrete unproductive route for a
zero-slot reduction: staying within the Laderman component and trying to
make a product vanish. Further rank-reduction searches must leave that
component or pursue a different decomposition mechanism. The output-tight
profile theorem gives a smaller, explicitly conditional branch of a
possible classification argument. What is still needed in that branch is
a proof that the five rank-two and seventeen rank-one outputs cannot
satisfy all the cross-block identities, or a construction showing they
can. Neither conclusion is asserted.

The repository's rules require a full resolution of the original field,
model, assumptions, and quantifiers for a solved status; a checked special
case does not suffice \cite{contributing}. This archive should therefore
be treated as partial research with exact supporting certificates, not as
a submission claiming that AC-02 is solved. No repository files were
modified.

\appendix
\section{The explicit Laderman identity used here}
\label{app:laderman}
This appendix uses the conventional 1-based matrix indices. Let
$A=(a_{ij})$, $B=(b_{ij})$, and define the following products. This is the
existing Laderman identity \cite{laderman}, not a new algorithm.
{\small
\begin{align*}
p_{1}&=(a_{11}+a_{12}+a_{13}-a_{21}-a_{22}-a_{32}-a_{33})(b_{22})\\
p_{2}&=(a_{11}-a_{21})(-b_{12}+b_{22})\\
p_{3}&=(a_{22})(-b_{11}+b_{12}+b_{21}-b_{22}-b_{23}-b_{31}+b_{33})\\
p_{4}&=(-a_{11}+a_{21}+a_{22})(b_{11}-b_{12}+b_{22})\\
p_{5}&=(a_{21}+a_{22})(-b_{11}+b_{12})\\
p_{6}&=(a_{11})(b_{11})\\
p_{7}&=(-a_{11}+a_{31}+a_{32})(b_{11}-b_{13}+b_{23})\\
p_{8}&=(-a_{11}+a_{31})(b_{13}-b_{23})\\
p_{9}&=(a_{31}+a_{32})(-b_{11}+b_{13})\\
p_{10}&=(a_{11}+a_{12}+a_{13}-a_{22}-a_{23}-a_{31}-a_{32})(b_{23})\\
p_{11}&=(a_{32})(-b_{11}+b_{13}+b_{21}-b_{22}-b_{23}-b_{31}+b_{32})\\
p_{12}&=(-a_{13}+a_{32}+a_{33})(b_{22}+b_{31}-b_{32})\\
p_{13}&=(a_{13}-a_{33})(b_{22}-b_{32})\\
p_{14}&=(a_{13})(b_{31})\\
p_{15}&=(a_{32}+a_{33})(-b_{31}+b_{32})\\
p_{16}&=(-a_{13}+a_{22}+a_{23})(b_{23}+b_{31}-b_{33})\\
p_{17}&=(a_{13}-a_{23})(b_{23}-b_{33})\\
p_{18}&=(a_{22}+a_{23})(-b_{31}+b_{33})\\
p_{19}&=(a_{12})(b_{21})\\
p_{20}&=(a_{23})(b_{32})\\
p_{21}&=(a_{21})(b_{13})\\
p_{22}&=(a_{31})(b_{12})\\
p_{23}&=(a_{33})(b_{33})
\end{align*}
}

The output entries are
\begin{align*}
c_{11}&=p_{6}+p_{14}+p_{19}\\
c_{12}&=p_{1}+p_{4}+p_{5}+p_{6}+p_{12}+p_{14}+p_{15}\\
c_{13}&=p_{6}+p_{7}+p_{9}+p_{10}+p_{14}+p_{16}+p_{18}\\
c_{21}&=p_{2}+p_{3}+p_{4}+p_{6}+p_{14}+p_{16}+p_{17}\\
c_{22}&=p_{2}+p_{4}+p_{5}+p_{6}+p_{20}\\
c_{23}&=p_{14}+p_{16}+p_{17}+p_{18}+p_{21}\\
c_{31}&=p_{6}+p_{7}+p_{8}+p_{11}+p_{12}+p_{13}+p_{14}\\
c_{32}&=p_{12}+p_{13}+p_{14}+p_{15}+p_{22}\\
c_{33}&=p_{6}+p_{7}+p_{8}+p_{9}+p_{23}
\end{align*}

The data ordering, support variety, signature slot labels, and all exact
certificates refer to this particular displayed ordering.

\section{Certificate format and audit map}
\label{app:cert}
All JSON label triples $(g,t,k)$ are 0-based: $g=0,1,2$ denotes $U,V,W$,
$t$ is the slot, and $k$ is the row-major coordinate. Restricted variables
are enumerated lexicographically over those triples where the baseline
entry is nonzero. A Brent polynomial has index $81a+9b+c$.

\begin{longtable}{p{0.39\textwidth}p{0.54\textwidth}}
\toprule
Certificate & What must be checked \\
\midrule\endhead
\file{*_forcing.json} & Each referenced equation has two nonzero-coefficient monomials; the designated known monomial is nonzero; the listed new variables are exactly the new factors forced. \\
\file{*_pairwise.json} & In each of three mode pairs, each selected row is supported on a unique remaining column and its two coefficients were forced nonzero. \\
\file{laderman_torus.json} & Initial and final selected minors have integer inverses; pivot formulas and identity free blocks hold; integer row combinations prove four cancellations; all 729 Laurent residuals vanish. \\
\file{laderman_gauge_slice.json} & The gauge exponent matrix is reconstructed; $ED=B$; the selected gauge minor is unimodular; the six remaining parameter indices and exponent columns agree. \\
\file{laderman_component.json} & $JH=0$ over the integers; the specified 545- and 76-dimensional minors are nonsingular modulo the verified prime. \\
\file{sun_tight.json} & Every output factorization and its rank is exact; $LR=RL=I_{27}$; all 529 cross-block identities hold. \\
\bottomrule
\end{longtable}

\textbf{No hidden oracle.} The generator may use SymPy to find inverse
matrices and decompositions, but the checker verifies the resulting
finite identities directly. The correctness argument for the component
theorem additionally uses the explicitly written algebraic-geometric
reasoning in Section \ref{sec:component}; a successful Python run alone
is not a formal verification of that reasoning.

\textbf{Integrity and reproducibility.} \file{MANIFEST.sha256} records the
files in the distributed archive. The environment versions and exact
run commands are in \file{environment.json} and \file{README.md}.
Historical source papers are cited rather than bundled. Numeric residuals
can vary slightly with numerical linear-algebra libraries, whereas the
exact verification assertions must have the same truth value.

\begin{thebibliography}{99}
\bibitem{ac02}
\emph{Open Problems in Numerical Linear Algebra}, AC-02,
``Exact bilinear rank of the $3\times3$ matrix product.'' Canonical
repository statement, last checked in the entry 10 September 2026;
accessed 14 September 2026.
\url{https://github.com/ajt60gaibb/OpenProblemsInNLA/tree/main/arithmetic-and-complexity/AC-02}

\bibitem{blaser}
M. Bl\"aser, \emph{Fast Matrix Multiplication}, Theory of Computing Library,
Graduate Surveys 5 (2013), 1--60. Introduction and Example 5.10.
DOI: 10.4086/toc.gs.2013.005.
\url{https://theoryofcomputing.org/articles/gs005/gs005.pdf}

\bibitem{laderman}
J. D. Laderman, ``A noncommutative algorithm for multiplying $3\times3$
matrices using 23 multiplications,'' Bulletin of the American Mathematical
Society 82 (1976), 126--128.
DOI: 10.1090/S0002-9904-1976-13988-2.

\bibitem{sun}
Y. Sun, \emph{An Exact 56-Addition, Rank-23 Scheme for General $3\times3$
Matrix Multiplication}, arXiv:2604.27645v1 (2026), Section 5.
\url{https://arxiv.org/html/2604.27645v1}

\bibitem{wang}
C. Wang, \emph{Automated Lower Bounds for Bilinear Complexity over Finite
Fields}, arXiv:2603.07280v11 (2026).
\url{https://arxiv.org/html/2603.07280v11}

\bibitem{almanli}
J. Alman and B. Li, \emph{Asymptotic Rank Speedup Theorems, Revisited},
arXiv:2605.21738v1 (2026), introduction.
\url{https://arxiv.org/html/2605.21738v1}

\bibitem{tichavsky}
P. Tichavsk\'y, \emph{Characterization of Decomposition of Matrix
Multiplication Tensors}, arXiv:2104.05323v1 (2021), Section IV and
Proposition 2.
\url{https://arxiv.org/html/2104.05323v1}

\bibitem{lbz}
X. Li, Y. Bao, and L. Zhang, \emph{On the local dimensions of solutions
of Brent equations}, arXiv:2303.09754v2 (2024).
\url{https://arxiv.org/html/2303.09754v2}

\bibitem{hks}
M. J. H. Heule, M. Kauers, and M. Seidl, \emph{New ways to multiply
$3\times3$-matrices}, arXiv:1905.10192v1 (2019), Section 7.
\url{https://arxiv.org/html/1905.10192v1}

\bibitem{contributing}
\emph{Open Problems in Numerical Linear Algebra}, contribution and
resolution requirements; accessed 14 September 2026.
\url{https://github.com/ajt60gaibb/OpenProblemsInNLA/blob/main/CONTRIBUTING.md}
\end{thebibliography}
\end{document}
